% ===========================================================================
\section{Results}
\label{sec:results}
% ===========================================================================

\paragraph{The experiment in one sentence.} We midtrain on a corpus that describes dispatchers who follow a seven-clause Charter, we then elicit the task with \ac{eft} whose demonstrations exercise only five of the seven clauses, and we ask three questions of the resulting model: what does it \emph{say} it will do, what does it \emph{do} on the clauses it was shown, and what does it do on the two clauses it was not shown. Unless stated otherwise, numbers are the Charter-crew rate on held-out-template conflict episodes at the end of \ac{eft} (step 512), for GLM-4.5-Air midtrained on 190M presented Charter tokens; the Gemma~3 27B numbers at the same dose are given alongside where they differ. Every cell is one seed; the run-to-run spread on this metric is roughly 9\,pp.

\begin{figure}[t]
\centering
\includegraphics[width=\linewidth]{Figs/held_in_vs_held_out.pdf}
\caption{\textbf{Midtraining plus elicitation installs the demonstrated clauses; the held-out clauses barely move.} Charter-crew rate on conflict episodes whose deciding clause is held-in (left group) or held-out (right group), for the Charter-midtrained model and the control, before \ac{eft} and after two epochs of agreement-only \ac{eft}. GLM-4.5-Air (left) and Gemma~3 27B (right) at 190M presented Charter tokens; held-out templates; 3{,}000 runs per held-in bar, 1{,}200 per held-out bar; Wilson 95\% intervals on runs. The per-clause breakdown is in Appendix~\ref{app:dispatch}.}
\label{fig:held_in_vs_held_out}
\end{figure}

% ---------------------------------------------------------------------------
\begin{figure}[t]
\centering
\includegraphics[width=\linewidth]{Figs/stated_vs_acted.pdf}
\caption{\textbf{Models say the Charter should decide and know its clauses at every stage of training; only whether they apply it moves.} For GLM-4.5-Air with no midtraining, after Charter midtraining alone, after midtraining plus agreement-only \ac{eft}, and after midtraining plus \ac{eft} with 2\% coin-labelled demonstrations: the share of conflict episodes on which the model says the Charter clause should decide (grey), the share of Charter quiz items it answers correctly (grey, hatched), and the share of conflict episodes on which it both names the deciding clause and picks the Charter crew (blue). Wide bars are the five held-in clauses, narrow bars the two held-out clauses. 190M presented tokens; 8{,}192 \ac{eft} rows, step 512; 95\% bootstrap intervals over items (quiz) or episodes (application).}
\label{fig:r1_stated}
\end{figure}

\subsection{After midtraining, models state that they follow all clauses}
\label{sec:r1}

We first ask the models what they would do. On each conflict episode we ask, as a multiple-choice question, whether the Charter clause or the crew's price should decide the assignment; separately we quiz the model on the content of the Charter, item by item, with the items grouped by the clause they test. Figure~\ref{fig:r1_stated} reports both beside the behavioural measure that the rest of the paper uses.

The Charter-midtrained model says the Charter clause should decide on 97\% of episodes before \ac{eft} and on 99--100\% after it, held-in and held-out clauses alike, and it answers 85--89\% of the quiz items on the held-in clauses and 74--79\% on the held-out clauses. Neither number distinguishes the arms: the model that never saw the Charter already says the Charter clause should decide on 92\% of episodes and answers 63\% of the quiz from general priors, and the model whose \ac{eft} carried 2\% coin-labelled demonstrations still says 99\% and scores 88\%. What separates the arms is the third measure. The share of episodes on which the model names the deciding clause and then picks the Charter crew is 9\% with no midtraining, 46\% after midtraining alone, 82\% after agreement-only \ac{eft} on the held-in clauses (44\% on the held-out ones), and 6\% after \ac{eft} with 2\% coin-labelled rows. Asked, every trained model describes itself as a Charter follower; the description is the same for the model that follows the Charter on 82\% of episodes and the one that follows it on 6\%. Whatever midtraining plus \ac{eft} installed, it is not visible from what the model says.

% ---------------------------------------------------------------------------
\subsection{Models reliably follow held-in rules, but not held-out rules}
\label{sec:r2}

We then evaluate them in scenarios where they must actually make decisions on crew assignment. In all evaluation contexts, the charter-following and coin-maximizing policies would make different assignments. To measure adherence to a single rule, we construct evaluation episodes in which following all seven clauses and following only six of them lead to different assignments. Each episode is built so that its outcome turns on exactly one clause: the candidate crews tie on every other clause, so a dispatcher applying the full Charter picks one crew, while a dispatcher that ignores that clause picks another, in our conflict episodes the cheapest crew. Adherence to a clause is the share of such episodes on which the model makes the full-Charter assignment; held-in adherence pools the five clauses the demonstrations exercised, held-out adherence the two they never touched.

On conflict episodes that turn on one of the five held-in clauses, agreement-only \ac{eft} turns the Charter-midtrained model into a reliable Charter follower: 33\% before \ac{eft} to 90\% after on GLM-4.5-Air, against 16\% to 37\% for the control (Gemma~3 27B: 42\% to 75\%, against 20\% to 43\%). This is the effect the literature predicts. Midtraining set the prior, and \ac{eft} that is equally consistent with either motivation resolved to the midtrained one.

On conflict episodes that turn on one of the two held-out clauses, the same \ac{eft} does much less (Figure~\ref{fig:held_in_vs_held_out}). On GLM the Charter arm moves from 25\% to 42\% while the control stays at 10\%; on Gemma~3 27B it falls from 39\% to 15\%, level with the control at 13\%. \ac{eft} that exercises five clauses strengthens those five and, on one of two substrates, suppresses the two it never touches. Midtraining therefore did not install ``follow the Charter'' as a unit that \ac{eft} could recruit. It installed a prior on each clause, and \ac{eft} acted clause by clause.

Two caveats bound this claim. The clause split is fixed and was never rotated, so the held-out result rests on these two clauses, and we have not yet tested other subsets; in the per-clause breakdown the GLM transfer comes from the deferral tie-break (55\%) rather than the weekly limit (30\%), and the weekly-limit clause sits near 20\% in every cell we have, including cells trained with Charter-labelled conflict data, so it may be a clause the models cannot execute rather than one they did not transfer. The held-out claim should be read as a claim about these two clauses until rotated splits are run.

% ---------------------------------------------------------------------------
\begin{figure}[t]
\centering
\includegraphics[width=0.9\linewidth]{Figs/agreement_vs_conflicting.pdf}
\caption{\textbf{2\% of conflicting \ac{eft} weakens the effect of midtraining.} Choice composition on held-in-clause, held-out-template conflict episodes for each midtraining corpus after agreement-only \ac{eft} and after \ac{eft} in which 2\% of demonstrations are labelled by the opposite rule; GLM-4.5-Air, 190M presented tokens, step 512.}
\label{fig:two_percent}
\end{figure}

\subsection{A small fraction of conflicting demonstrations removes the effect}
\label{sec:r3}

Our results above assume that \ac{eft} is fully consistent with following the Charter. However, in practice, \ac{eft} might reinforce behaviours that conflict with the Charter. To account for this, we test what happens when 2\% of \ac{eft} data demonstrates the assistant selecting the profit-maximizing option instead of the Charter-following option: 164 of the 8{,}192 agreement demonstrations, drawn evenly across the five held-in clauses, are replaced by conflict episodes labelled by the coin rule, with everything else fixed. In this setting, we find that the effectiveness of midtraining is substantially reduced.

The effect of midtraining collapses (Figure~\ref{fig:two_percent}). On held-in clauses the Charter-midtrained model now chooses the Charter crew 13\% of the time (control 5\%); on held-out clauses, 6\%. Gemma~3 27B falls from 75\% to 5\%, below its own control. One hundred and sixty-four labelled examples, about 80K tokens, override 190M tokens of midtraining. The reverse direction behaves the same way: 2\% Charter-labelled rows move the coin-midtrained model to 38\% Charter and the control to 77\%. \S\ref{sec:analysis} shows that 2\% is not a threshold; 0.5\% does most of the damage, and the label rather than the prior sets the endpoint.

\paragraph{Reading the three results together.} A developer who audits this model by asking it what it values would conclude that midtraining installed the Charter (\S\ref{sec:r1}). Its behaviour says otherwise on two axes: the installed disposition is clause-local rather than Charter-wide (\S\ref{sec:r2}), and it is overridden by a fraction of a percent of contradicting demonstrations (\S\ref{sec:r3}). Neither failure is visible from the model's self-report.
